\centerline{\bf \Huge Обрыгонометрия}

\begin{flushright}
        \scriptsize{Эпизод 22. Оставаться человеком}
\end{flushright}

\problem{1}
Число $x$ таково, что обе суммы $S=\sin 64 x+\sin 65 x$ и $C=\cos 64 x+\cos 65 x$ - рациональные числа. Докажите, что в одной из этих сумм оба слагаемых рациональны.
%https://3.shkolkovo.online/catalog/7853

\problem{2}
Даны вещественные числа $x_1, \ldots, x_n$. Найдите максимальное значение выражения

$$
A=\frac{\sin x_1+\ldots+\sin x_n}{\sqrt{\operatorname{tg}^2 x_1+\ldots+\operatorname{tg}^2 x_n+n}} .
$$

\problem{3}
Найдите угол $\alpha$, если известно, что $0<\alpha<90^{\circ}$ и

$$
\operatorname{tg} \alpha=\frac{\left(1+\operatorname{tg} 2^{\circ}\right)\left(1+\operatorname{tg} 5^{\circ}\right)-2}{\left(1-\operatorname{tg} 2^{\circ}\right)\left(1-\operatorname{tg} 5^{\circ}\right)-2}
$$

\problem{4}
Существует ли такое положительное число $a$, что при всех действительных $x$ верно неравенство

$$
|\cos x|+|\cos a x|>\sin x+\sin a x ?
$$

\problem{5}
Решите неравенство

$$
\arcsin (\sin |x|) \geqslant \arccos |\cos 3 x|
$$

\problem{6}
Докажите, что при $k>10$ в произведении

$$
f(x)=\cos (x) \cos (2 x) \cos (3 x) \ldots \cos \left(2^k x\right)
$$


можно заменить один $\cos$ на $\sin$ так, что получится функция $f_1(x)$, удовлетворяющая при всех действительных $x$ неравенству $\left|f_1(x)\right| \leqslant \frac{3}{2^{k+1}}$.

\problem{7}
Известно, что для некоторых $x$ и $y$ суммы $\sin x+\cos y$ и $\sin y+\cos x$ - положительные рациональные числа. Докажите, что найдутся такие натуральные числа $m$ и $n$, что $m \sin x+n \cos x$ - натуральное число.